\documentclass[12pt,a4paper]{article}
\usepackage[utf8]{inputenc}
\usepackage[T1]{fontenc}
\usepackage[english,romanian]{babel}
\usepackage{graphicx}
\usepackage{subcaption}
\usepackage{amsmath,amssymb}
\usepackage{fancyhdr}
\usepackage[a4paper,margin=2.5cm]{geometry}
\usepackage{dirtytalk}
\usepackage{csquotes}
\usepackage[hidelinks]{hyperref}
\usepackage{listings}
\usepackage{xcolor}
\usepackage{array}
\usepackage{booktabs}

\lstset{
    language=Python,
    basicstyle=\ttfamily\small,
    numbers=left,
    numberstyle=\tiny,
    frame=single,
    breaklines=true,
    showstringspaces=false,
    captionpos=b
}
% listings nu înțelege direct diacriticele UTF-8, de aceea le traducem explicit.
\lstset{
    literate=
        {Ă}{{\u{A}}}1 {ă}{{\u{a}}}1
        {Â}{{\^A}}1 {â}{{\^a}}1
        {Î}{{\^I}}1 {î}{{\^i}}1
        {Ș}{{\textcommabelow{S}}}1 {ș}{{\textcommabelow{s}}}1
        {Ț}{{\textcommabelow{T}}}1 {ț}{{\textcommabelow{t}}}1
}
\lstdefinestyle{terminal}{
    language={},
    numbers=none,
    backgroundcolor=\color{black!5},
    basicstyle=\ttfamily\footnotesize
}
\setlength{\parindent}{1.25cm}
\setlength{\parskip}{0.25em}

\begin{document}
\begin{titlepage}
\begin{center}
    % Înlocuiți calea dacă logoul este stocat în altă parte.
    \includegraphics[width=0.3\textwidth]{lab1/image .png}\\[0.2cm]
    \textbf{MINISTERUL EDUCAȚIEI, CULTURII ȘI CERCETĂRII\\
    AL REPUBLICII MOLDOVA}\\[0.3cm]
    \textbf{Universitatea Tehnică a Moldovei\\
    Facultatea Calculatoare, Informatică și Microelectronică\\
    Departamentul Inginerie Software și Automatică}\\[2cm]
    \Huge\textbf{Raport}\\[0.5cm]
    \Large Lucrarea de laborator nr. 3\\[0.3cm]
    \large\textbf{la Criptografie și Securitate}\\[0.8cm]
    \textbf{Ciur Alexandrina, FAF-241}\\[2.5cm]
    \begin{flushright}
        \textit{Verificat de:}\\
        \textbf{Zaica M.}\\
        DISA, FCIM, UTM
    \end{flushright}
    \vfill
    Chișinău -- 2026
\end{center}
\end{titlepage}

\section*{Tema: Cifruri polialfabetice. Algoritmul Playfair}

\section{Scopul lucrării}

Studierea cifrurilor de substituție polialfabetică și implementarea algoritmului
de criptare Playfair pentru mesaje scrise în limba română (alfabet de 31 de litere),
cu posibilitatea de criptare și decriptare a textului.

\section{Sarcina lucrării}

Conform registrului grupei, studentul are numărul 7 (număr impar), prin urmare
a fost realizată \textbf{Sarcina 3.1}:

\begin{quote}
\itshape
De implementat algoritmul Playfair în unul din limbajele de programare pentru
mesaje în limba română (31 de litere). Valorile caracterelor textului sunt cuprinse
între `A' și `Z', `a' și `z' și nu sunt permise alte valori. În cazul în care
utilizatorul introduce alte valori -- i se va sugera diapazonul corect al
caracterelor. Lungimea cheii nu trebuie să fie mai mică de 7. Utilizatorul va putea
alege operația -- criptare sau decriptare, va putea introduce cheia, mesajul sau
criptograma și va obține criptograma sau mesajul decriptat. Faza finală de adăugare
a spațiilor noi, în funcție de limba folosită și de logica mesajului -- se va face
manual.
\end{quote}

\section{Considerații teoretice}

\subsection{Cifrurile polialfabetice}

Slăbiciunea cifrurilor monoalfabetice (de exemplu, cifrul Cezar) constă în faptul
că fiecare literă a textului clar este înlocuită mereu cu aceeași literă a
criptogramei. Astfel, distribuția frecvențelor literelor din criptogramă repetă
distribuția din limba textului clar, iar cifrul poate fi spart prin analiza
frecvențelor.

La cifrurile \textbf{polialfabetice}, substituția unui caracter variază în funcție
de diverși parametri (poziția în text, contextul, literele vecine etc.). Aceeași
literă a textului clar poate fi criptată în moduri diferite, iar distribuția
frecvențelor devine mai uniformă și oferă mai puține informații criptanalistului.
Exemple clasice sunt cifrul Vigenère și cifrul Playfair.

\subsection{Scurt istoric al algoritmului Playfair}

Algoritmul a fost inventat de Charles Wheatstone și descris pentru prima dată la
26 martie 1854, însă poartă numele baronului Lyon Playfair, care a convins guvernul
britanic să-l adopte pentru uz oficial. A fost folosit de armata britanică în
războiul cu burii, iar variante modificate au fost folosite în Primul și al Doilea
Război Mondial. Astăzi algoritmul este considerat învechit, dar este unul dintre
primii algoritmi care criptează \textbf{blocuri} de câte două litere (digrame),
idee folosită și în cifrurile bloc moderne.

\subsection{Descrierea algoritmului Playfair}

Criptarea Playfair constă din trei etape:
\begin{enumerate}
    \item pregătirea textului ce urmează a fi criptat;
    \item construirea matricei de criptare pe baza cheii;
    \item criptarea textului pe perechi de litere.
\end{enumerate}

\paragraph{1. Pregătirea textului.} Textul se scrie cu majuscule, fără spații și
semne de punctuație, apoi se împarte în perechi de câte două litere. Dacă într-o
pereche ar apărea două litere identice (de exemplu \texttt{LL}), între ele se
inserează o literă de umplutură (de regulă una rar folosită, cum ar fi \texttt{X}).
Dacă la final rămâne o singură literă, i se adaugă o literă de umplutură.

\paragraph{2. Construirea matricei.} Din cheie se elimină literele care se repetă,
păstrându-se doar prima apariție. Matricea se completează pe rânduri, începând din
colțul stânga-sus, mai întâi cu literele cheii, apoi cu restul literelor
alfabetului, în ordine alfabetică. Fiecare literă apare în matrice o singură dată.

\paragraph{3. Criptarea perechilor.} Fie perechea de litere aflate pe pozițiile
$(r_1, c_1)$ și $(r_2, c_2)$ într-o matrice cu $R$ rânduri și $C$ coloane:
\begin{itemize}
    \item \textbf{aceeași linie} ($r_1 = r_2$) -- fiecare literă se înlocuiește cu
    cea din dreapta ei (ultima trece la prima):
    \[
        (r, c) \rightarrow \bigl(r,\ (c + 1) \bmod C\bigr);
    \]
    \item \textbf{aceeași coloană} ($c_1 = c_2$) -- fiecare literă se înlocuiește
    cu cea de dedesubt (ultima trece la prima):
    \[
        (r, c) \rightarrow \bigl((r + 1) \bmod R,\ c\bigr);
    \]
    \item \textbf{linii și coloane diferite} -- literele formează colțurile unui
    dreptunghi și fiecare se înlocuiește cu litera de pe aceeași linie, dar de pe
    coloana celeilalte litere:
    \[
        (r_1, c_1),\ (r_2, c_2) \rightarrow (r_1, c_2),\ (r_2, c_1).
    \]
\end{itemize}

\paragraph{Decriptarea} folosește aceeași matrice și aceleași reguli, dar în sens
invers: pe aceeași linie se ia litera din \textbf{stânga}
($(c - 1) \bmod C$), pe aceeași coloană -- litera de \textbf{deasupra}
($(r - 1) \bmod R$), iar regula dreptunghiului rămâne neschimbată. După decriptare,
literele de umplutură se elimină și spațiile se adaugă manual, în funcție de
sensul mesajului.

\subsection{Adaptarea pentru alfabetul român}

Alfabetul român are 31 de litere:
\begin{center}
\texttt{A Ă Â B C D E F G H I Î J K L M N O P Q R S Ș T Ț U V W X Y Z}
\end{center}

Numărul 31 este prim, deci literele nu pot fi aranjate exact într-o matrice
dreptunghiulară fără celule libere. Conform indicațiilor lucrării, se poate elimina
litera cel mai rar întâlnită în limba respectivă, iar la decriptare ea se restabilește
intuitiv. În implementare au fost luate următoarele decizii:
\begin{itemize}
    \item litera \textbf{Q} (cea mai rară în limba română) este înlocuită cu
    \textbf{K} atât în cheie, cât și în mesaj;
    \item rămân $31 - 1 = 30$ de litere, care formează o matrice de
    \textbf{6 linii $\times$ 5 coloane};
    \item litera de umplutură este \textbf{X}; dacă litera dublată este chiar
    \texttt{X}, se folosește \textbf{Y};
    \item lungimea minimă a cheii este de 7 litere (fără spații);
    \item se acceptă doar literele \texttt{A--Z}, \texttt{a--z},
    \texttt{Ă, Â, Î, Ș, Ț} (și minusculele lor) și spații.
\end{itemize}

Numărul de matrice posibile este $30! \approx 2{,}65 \cdot 10^{32}$, ceea ce face
imposibilă căutarea exhaustivă manuală a cheii.

\subsection{Exemplu de criptare manuală}

Pentru cheia \enquote{Primăvara frumoasă}, după eliminarea spațiilor și a literelor
repetate, obținem \texttt{PRIMĂVAFUOS}. Matricea de criptare este:

\begin{center}
\renewcommand{\arraystretch}{1.3}
\begin{tabular}{|c|c|c|c|c|}
\hline
P & R & I & M & Ă \\ \hline
V & A & F & U & O \\ \hline
S & Â & B & C & D \\ \hline
E & G & H & Î & J \\ \hline
K & L & N & Ș & T \\ \hline
Ț & W & X & Y & Z \\ \hline
\end{tabular}
\end{center}

Exemple de aplicare a fiecărei reguli:

\begin{center}
\renewcommand{\arraystretch}{1.3}
\begin{tabular}{@{}cccl@{}}
\toprule
\textbf{Pereche} & \textbf{Poziții} & \textbf{Rezultat} & \textbf{Regula aplicată} \\
\midrule
PI & (0,0), (0,2) & RM & aceeași linie $\rightarrow$ litera din dreapta \\
IN & (0,2), (4,2) & FX & aceeași coloană $\rightarrow$ litera de dedesubt \\
MÂ & (0,3), (2,1) & RC & dreptunghi $\rightarrow$ se schimbă coloanele \\
\bottomrule
\end{tabular}
\end{center}

Mesajul \enquote{Accesul la date secrete} se pregătește astfel: se scot spațiile
(\texttt{ACCESULLADATESECRETE}), se împarte în perechi, între cele două litere
\texttt{L} din aceeași pereche se inserează \texttt{X}, iar la final se adaugă
\texttt{X}, deoarece rămâne o literă fără pereche:
\begin{center}
\texttt{AC CE SU LX LA DA TE SE CR ET EX}
\end{center}
Criptograma obținută este \texttt{UÂSÎCVNWWÂÂOKJEKÂMJKHȚ}.

\section{Implementarea}

Programul a fost realizat în limbajul \textbf{Python 3} (fișierul
\texttt{lab3/lab3.py}). Mai jos sunt prezentate părțile principale ale codului.

\subsection{Alfabetul și constantele}

\begin{lstlisting}[caption={Alfabetul român și parametrii matricei}]
ALPHABET = (
    "A", "Ă", "Â", "B", "C", "D", "E", "F", "G", "H", "I", "Î", "J", "K",
    "L", "M", "N", "O", "P", "Q", "R", "S", "Ș", "T", "Ț", "U", "V", "W",
    "X", "Y", "Z",
)
MIN_KEY_LENGTH = 7
PLAYFAIR_ROWS = 6
PLAYFAIR_COLUMNS = 5
# 31 de litere nu umplu o matrice dreptunghiulara, deci Q (cea mai rara litera
# in limba romana) este inlocuita cu K si raman 30 de litere = 6 x 5.
REPLACED_LETTER = "Q"
REPLACEMENT = "K"
MATRIX_ALPHABET = tuple(
    character for character in ALPHABET if character != REPLACED_LETTER
)
FILLER = "X"
ALTERNATE_FILLER = "Y"
\end{lstlisting}

Tupla \texttt{ALPHABET} conține toate cele 31 de litere ale alfabetului român.
\texttt{MATRIX\_ALPHABET} este același alfabet fără litera \texttt{Q}, adică cele
30 de litere care se plasează în matricea $6 \times 5$.

\subsection{Validarea textului introdus}

\begin{lstlisting}[caption={Normalizarea și validarea textului}]
def validate_text(raw):
    """Validate and compact a message or key, returning (text, invalid_character)."""
    text = normalize(raw).replace(" ", "")
    for character in text:
        if character not in ALPHABET:
            return None, character
    return text.replace(REPLACED_LETTER, REPLACEMENT), None
\end{lstlisting}

Funcția transformă textul în majuscule, elimină spațiile și verifică fiecare
caracter. Dacă este găsit un caracter nepermis (cifră, semn de punctuație etc.),
acesta este returnat, iar programul îi afișează utilizatorului diapazonul corect
al caracterelor. La final, litera \texttt{Q} este înlocuită cu \texttt{K}.

\subsection{Construirea matricei Playfair}

\begin{lstlisting}[caption={Eliminarea literelor repetate din cheie și construirea matricei}]
def unique_key(key):
    """Remove repeated letters from a Playfair key, preserving first occurrence."""
    return "".join(dict.fromkeys(key))


def build_playfair_matrix(key):
    """Build a 6x5 matrix containing every matrix alphabet character once."""
    key = unique_key(key)
    sequence = key + "".join(
        character for character in MATRIX_ALPHABET if character not in key
    )
    return [
        sequence[row * PLAYFAIR_COLUMNS:(row + 1) * PLAYFAIR_COLUMNS]
        for row in range(PLAYFAIR_ROWS)
    ]
\end{lstlisting}

\texttt{dict.fromkeys(key)} păstrează doar prima apariție a fiecărei litere din cheie.
La cheie se adaugă literele alfabetului care nu apar în ea, iar șirul de 30 de
litere obținut este tăiat în 6 rânduri a câte 5 litere.

\subsection{Pregătirea mesajului}

\begin{lstlisting}[caption={Împărțirea în perechi și inserarea literelor de umplutură}]
def _filler_for(character):
    return ALTERNATE_FILLER if character == FILLER else FILLER


def prepare_playfair_text(text):
    """Split text into digraphs, separating repeated letters and padding the end."""
    pairs = []
    index = 0
    while index < len(text):
        first = text[index]
        second = text[index + 1] if index + 1 < len(text) else None
        if second is None or first == second:
            pairs.append(first + _filler_for(first))
            index += 1
        else:
            pairs.append(first + second)
            index += 2
    return "".join(pairs)
\end{lstlisting}

Dacă perechea ar conține două litere identice sau dacă a rămas o singură literă la
sfârșit, se adaugă litera de umplutură, iar a doua literă trece în perechea următoare.

\subsection{Criptarea și decriptarea}

\begin{lstlisting}[caption={Aplicarea celor trei reguli Playfair}]
def playfair(text, key, decrypt=False):
    """Encrypt or decrypt a prepared Playfair message."""
    if len(text) % 2:
        raise ValueError("Textul Playfair trebuie sa contina un numar par de litere.")
    matrix = build_playfair_matrix(key)
    positions = _playfair_positions(matrix)
    result = []
    step = -1 if decrypt else 1

    for index in range(0, len(text), 2):
        first, second = text[index:index + 2]
        row1, column1 = positions[first]
        row2, column2 = positions[second]
        if row1 == row2:
            result.extend((
                matrix[row1][(column1 + step) % PLAYFAIR_COLUMNS],
                matrix[row2][(column2 + step) % PLAYFAIR_COLUMNS],
            ))
        elif column1 == column2:
            result.extend((
                matrix[(row1 + step) % PLAYFAIR_ROWS][column1],
                matrix[(row2 + step) % PLAYFAIR_ROWS][column2],
            ))
        else:
            result.extend((matrix[row1][column2], matrix[row2][column1]))
    return "".join(result)
\end{lstlisting}

Aceeași funcție realizează și criptarea, și decriptarea. Diferența este dată de
variabila \texttt{step}: la criptare este $+1$ (deplasare la dreapta/în jos), iar la
decriptare este $-1$ (deplasare la stânga/în sus). Operația modulo (\texttt{\%})
asigură trecerea de la ultima literă a liniei sau a coloanei la prima.

\subsection{Citirea datelor de la utilizator}

\begin{lstlisting}[caption={Citirea mesajului cu verificarea erorilor}]
def _read_message(prompt, decrypt=False):
    while True:
        message, invalid = validate_text(input(prompt))
        if invalid is not None:
            print(f"Caracter invalid: {invalid!r}. Folositi doar: {ALLOWED_HINT}.")
        elif not message:
            print("Textul nu poate fi gol.")
        elif decrypt and len(message) % 2:
            print("Criptograma trebuie sa aiba un numar par de litere.")
        else:
            return message
\end{lstlisting}

Programul cere datele din nou până când acestea sunt corecte: caracterele trebuie
să fie permise, textul nu poate fi gol, iar criptograma trebuie să aibă un număr par
de litere. Cheia este verificată analog, inclusiv lungimea minimă de 7 litere.

\section{Rezultatele execuției}

Programul se lansează din terminal cu comanda:
\begin{lstlisting}[style=terminal]
python lab3/lab3.py
\end{lstlisting}

\subsection{Criptarea mesajului}

Utilizatorul alege opțiunea \texttt{1}, introduce cheia și mesajul. Programul
afișează matricea, mesajul pregătit (împărțit în perechi) și criptograma.

\begin{lstlisting}[style=terminal, caption={Criptarea mesajului \enquote{Accesul la date secrete}}]
=== Cifrul Playfair (alfabetul romanesc) ===
1. Criptare
2. Decriptare
0. Iesire
Alege (0-2): 1
Cheia (minimum 7 litere): Primăvara frumoasă
Mesajul: Accesul la date secrete
Matricea Playfair:
P R I M Ă
V A F U O
S Â B C D
E G H Î J
K L N Ș T
Ț W X Y Z
Mesaj pregatit: AC CE SU LX LA DA TE SE CR ET EX
Criptograma: UÂSÎCVNWWÂÂOKJEKÂMJKHȚ
\end{lstlisting}

% Pentru a adăuga o captură de ecran, încărcați imaginea în Overleaf și
% decomentați blocul de mai jos (schimbați numele fișierului dacă e altul).
% \begin{figure}[h!]
%     \centering
%     \includegraphics[width=0.9\textwidth]{lab3/criptare.png}
%     \caption{Criptarea mesajului în terminal}
% \end{figure}

\subsection{Decriptarea criptogramei}

Utilizatorul alege opțiunea \texttt{2}, introduce \textbf{aceeași cheie} și
criptograma obținută anterior.

\begin{lstlisting}[style=terminal, caption={Decriptarea criptogramei}]
Alege (0-2): 2
Cheia (minimum 7 litere): Primăvara frumoasă
Criptograma: UÂSÎCVNWWÂÂOKJEKÂMJKHȚ
Matricea Playfair:
P R I M Ă
V A F U O
S Â B C D
E G H Î J
K L N Ș T
Ț W X Y Z
Mesaj decriptat (pe perechi): AC CE SU LX LA DA TE SE CR ET EX
Mesaj decriptat: ACCESULXLADATESECRETEX
Eliminati manual literele X/Y adaugate si puneti spatiile intre cuvinte.
\end{lstlisting}

% \begin{figure}[h!]
%     \centering
%     \includegraphics[width=0.9\textwidth]{lab3/decriptare.png}
%     \caption{Decriptarea criptogramei în terminal}
% \end{figure}

Eliminând manual literele de umplutură \texttt{X} (cea dintre cele două
\texttt{L} și cea de la final) și adăugând spațiile, obținem mesajul inițial:
\begin{center}
\texttt{ACCESUL LA DATE SECRETE}
\end{center}

\subsection{Tratarea datelor incorecte}

Dacă utilizatorul introduce o opțiune inexistentă, o cheie prea scurtă, caractere
nepermise sau o criptogramă cu număr impar de litere, programul afișează un mesaj
explicativ și cere datele din nou.

\begin{lstlisting}[style=terminal, caption={Mesaje de eroare la introducerea datelor greșite}]
Alege (0-2): 5
Optiune invalida. Alegeti un numar intre 0 si 2.

Alege (0-2): 1
Cheia (minimum 7 litere): cheie
Cheia trebuie sa aiba cel putin 7 litere.
Cheia (minimum 7 litere): Primăvara frumoasă
Mesajul: Parola: 1234
Caracter invalid: ':'. Folositi doar: A-Z, a-z, Ă, Â, Î, Ș, Ț (ă, â, î, ș, ț) si spatii.
...
Alege (0-2): 2
Cheia (minimum 7 litere): Primăvara frumoasă
Criptograma: UÂSÎC
Criptograma trebuie sa aiba un numar par de litere.
\end{lstlisting}

\section{Concluzii}

În cadrul acestei lucrări de laborator a fost studiat principiul de funcționare al
cifrurilor polialfabetice și a fost implementat algoritmul Playfair pentru
alfabetul român în limbajul Python. Programul permite alegerea operației (criptare
sau decriptare), introducerea cheii și a mesajului sau a criptogramei, verifică
corectitudinea datelor și sugerează diapazonul corect al caracterelor în caz de
eroare.

Deoarece alfabetul român are 31 de litere, iar 31 este un număr prim, a fost
necesară adaptarea algoritmului: litera cea mai rară, \texttt{Q}, a fost înlocuită
cu \texttt{K}, obținându-se o matrice de $6 \times 5$. Am observat că, spre deosebire
de cifrurile monoalfabetice, la Playfair aceeași literă se poate cripta diferit, în
funcție de litera cu care formează pereche, ceea ce ascunde frecvențele literelor
individuale. Totuși, algoritmul păstrează frecvențele perechilor de litere și poate
fi spart prin analiza digramelor, de aceea astăzi are doar valoare didactică.

De asemenea, am constatat că decriptarea reușește doar cu exact aceeași cheie ca la
criptare: o cheie diferită generează o altă matrice, iar rezultatul este un text
fără sens. Pasul final -- eliminarea literelor de umplutură și adăugarea spațiilor --
se face manual, deoarece programul nu poate deosebi o literă \texttt{X} adăugată de
una care face parte din mesaj (de exemplu, în cuvântul \enquote{examen}).

Codul sursă al lucrării este disponibil în repozitoriul GitHub:\\
\url{https://github.com/alecsyyyk/Cryptography-and-Security}

\end{document}
